\section{Evaluation}
\label{sec:evaluation}

\subsection{Experimental Setup}
\begin{itemize}
    \item Report the pretrained checkpoint, encoder optimization scope, input
    construction, optimizer, learning-rate schedule, batch construction,
    maximum epochs, and early stopping.
    \item Report the three seeds used for the literature-facing joint-hierarchy
    benchmark and the single seed used for the prospective clean controls.
    \item Define all metric abbreviations at first use.
    \item \textit{Configuration placeholder: shared and stream-specific
    experimental settings.}
\end{itemize}

\subsection{Two Evidence Streams and Selection Protocols}
\begin{itemize}
    \item Define the literature-facing benchmark stream: task-compatible
    published ICBHI rows, verified zero-target transfer rows, and the
    three-seed joint-hierarchy result.
    \item Disclose that the three-seed joint-hierarchy checkpoints are selected
    by ICBHI official-test Score under the PAFA-style benchmark protocol;
    SPRSound inter is evaluated after each checkpoint is selected.
    \item Define the prospective clean stream: group-safe internal calibration
    and selection, validation-only checkpoint and threshold choice, and one
    official-test evaluation after freezing.
    \item Use one preregistered joint validation criterion for the two joint
    clean controls while keeping native task metrics separate.
    \item Do not combine or relabel the benchmark and prospective clean rows.
\end{itemize}

\subsection{Literature-Facing Benchmark and Transfer Motivation}
\begin{itemize}
    \item Include task-compatible published ICBHI references with their task,
    unit, split, seed count, and selection disclosure.
    \item Include verified zero-target SPRSound transfer of fixed
    ICBHI-specialized PAFA, SG-SCL, and Patch-Mix checkpoints as motivation.
    \item Identify the all-Normal floor used to interpret zero-target transfer.
    \item Include the three-seed ICBHI and SPRSound native-task results for the
    proposed joint hierarchy.
\end{itemize}

\begin{table*}[t]
    \centering
    \fbox{\parbox[c][1.05in][c]{0.96\textwidth}{
        \centering
        \textit{Table I placeholder: test-selected task-compatible ICBHI
        references; fixed-checkpoint zero-target SPRSound transfer motivation;
        and the three-seed joint-hierarchy ICBHI/SPRSound benchmark with
        selection disclosure.}
    }}
    \caption{\textit{Table I caption placeholder: literature-facing benchmark
    and zero-target transfer motivation. Keep task, unit, split, seed count,
    metric, and model-selection source explicit.}}
    \label{tab:benchmark_results}
\end{table*}

\subsection{Prospective Clean Controls}
\begin{itemize}
    \item Compare joint eligibility-aware hierarchy with joint independent
    ICBHI-flat4 and SPRSound-binary heads.
    \item Compare the joint hierarchy with ICBHI-only hierarchy training and a
    frozen-checkpoint SPRSound transfer.
    \item Compare the joint hierarchy with SPRSound-only hierarchy training and
    a frozen-checkpoint ICBHI transfer.
    \item Report seed42 validation-selected native metrics and the exact
    checkpoint-selection criterion for every row.
    \item Keep the hierarchy-versus-independent-head claim conditional until
    all four rows are complete.
\end{itemize}

\begin{table*}[t]
    \centering
    \fbox{\parbox[c][0.95in][c]{0.96\textwidth}{
        \centering
        \textit{Table II placeholder: prospective clean seed42 joint
        hierarchy, joint independent heads, ICBHI-only hierarchy, and
        SPRSound-only hierarchy under one frozen validation protocol.}
    }}
    \caption{\textit{Table II caption placeholder: validation-selected clean
    controls. Report separate ICBHI and SPRSound native metrics; do not compare
    them through a pooled score.}}
    \label{tab:clean_controls}
\end{table*}

\subsection{Specificity, Sensitivity, and Per-Class Errors}
\begin{itemize}
    \item Report ICBHI specificity and sensitivity separately from aggregate
    Score.
    \item Report Normal, Crackle, Wheeze, and Both recall and class support.
    \item Identify the dominant confusion directions for the joint hierarchy
    and matched clean controls.
    \item Compare SPRSound Normal and Adventitious behavior using its native
    binary metrics.
    \item \textit{Analysis placeholder: determine whether retention or loss is
    concentrated in a dataset, hierarchy level, or low-support class.}
\end{itemize}

\subsection{Supporting Diagnostics and Discussion}
\begin{itemize}
    \item \textit{Bounded supporting-diagnostic placeholder: summarize fixed
    joint-hierarchy checkpoint behavior on HF Lung attributes and patient-grouped
    KAUH without presenting either as a native reproduction, main table, or
    principal contribution.}
    \item Discuss dataset-dependent retention and trade-offs using Figure~1,
    Table~I, Table~II, and the per-class analysis.
    \item Separate observed associations from causal explanations.
\end{itemize}
